\section{Weekly Summary}

\begin{tcolorbox}[colback=yellow!10,colframe=black!50,title=AI-Generated Content]
\noindent The summary below was written by Claude (Anthropic) from that week's regression outputs and series descriptions. It has not been reviewed or fact-checked by a human, and should be read as an automated, informal recap -- not analysis.
\end{tcolorbox}

Welcome back to the weekly ritual where we shove two economic time series that have never met into a blender and ask the resulting smoothie to explain itself. As always, the machine picked these pairs at random, nobody hypothesized anything, and any relationship you see is best understood as a cosmic coincidence wearing a lab coat. Let's marvel at the mundane together.

Monday, August 31, opened with a real tearjerker: a discontinued SBA loan series that gave up on life in 2017 got paired with the Ft. Lauderdale house price index, which has been faithfully tracking Florida real estate since disco was relevant. The raw regression produced an R-squared of 0.05 and a p-value of 0.33, which in technical terms means "the model shrugged." Then the detrended version nudged things ever so slightly upward to an R-squared of 0.10 with a p-value of 0.18 -- still nowhere near significant, but slightly less pathetic. So congratulations to nobody: this pairing achieved the rare feat of being uninteresting in both incarnations. Truly a series that quit its job in 2017 has nothing meaningful to say about Florida condos. Who could have foreseen it.

September 2 is where the plot thickens, and by "thickens" I mean "thins dramatically upon inspection." Aggregate weekly hours in durable goods manufacturing got hitched to demand deposits from foreign official institutions (also discontinued, because this project apparently runs a hospice for retired data series). The raw regression came back significant -- p-value of 0.006, hold the applause -- suggesting a genuine relationship to the untrained and hopeful eye. But then we detrended it, and the whole thing evaporated so fast it left skid marks: R-squared collapsed to 0.0003 and the p-value ballooned to 0.82. This, dear reader, is the textbook trend mirage this project exists to mock. Two lines drifting in the same general direction over decades, mistaken for romance. Strip out the shared drift and there's absolutely nothing there. Beautiful. Chef's kiss.

Finally, September 3 gave us imports of non-manufactured goods to the U.S. Virgin Islands versus Dutch services business confidence, a pairing so geographically and thematically unhinged I'm surprised the software didn't file a complaint. The raw regression was a dud (p-value 0.59), but the detrended version actually did something: p-value of 0.014, technically significant, with a modest R-squared of 0.027. So here's the fun twist -- it's the only pairing this week that got stronger after detrending rather than dying, which by the wretched standards of this project makes it the "notable" result. Do I believe Virgin Islands cargo manifests are secretly wired to the mood of Dutch service-sector managers? Absolutely not, and neither should you; an R-squared of 0.027 explains roughly nothing and a random search across enough series will cough up a 0.014 p-value eventually just to keep things spicy. But it holds up better than a shared trend, so it earns a polite golf clap before we forget it forever. Same time next week, same disclaimer: nothing here means anything, and that's the whole point.
